
% ════════════════════════════════════════════
% 模型建立（5.X.3.模型建立.tex）写作规则 —— 模板内嵌，AI 先读本文件再写
% ════════════════════════════════════════════
% 【定位】本块为每个问题的「模型建立」（三级子块），放在模型准备之后，用于建立本问的数学模型并展示中间结果。
% 【结构】开头段引文 → 分题型建模 → 结尾段总结。
% 【开头段（引文）】引出下面具体步骤讨论，说明本问模型的整体思路。
% 【题型分支（按题目主要类型选择；组合题型看主要类型，不需全部选）】：
%   · 评价：① 选指标；② 熵权法确定权重【中间结果需展示】；③ 任选一个评价方法进行运算。
%   · 分类：① 特征工程（选指标）【中间结果需展示】；② 选定模型（逻辑回归、Xgboost）；③ 模型评估（讲如何控制输出）。
%   · 预测：① 特征工程（选指标）【中间结果需展示】；② 选定模型（Xgboost）；③ 模型评估（讲如何控制输出）。
%   · 优化：① 建立目标函数；② 设定约束条件；③ 选定求解算法【开头要展示整个优化模型公式：目标函数+约束条件，然后再结合问题讨论】。
%   · 机理分析：AI 合理规划步骤，3-4 步即可，避免为凑公式强行堆砌。
% 【公式】每步至少 6 个公式，公式间用文字串联推导逻辑，每个公式后跟变量说明段。
% 【结尾段（总结）】总结本问建模成果，引出模型求解。
% 【表格】同 4.符号说明 的 longtable 规范（列宽总和=1.04-0.04*N，表头列名不加粗）。
% ════════════════════════════════════════════
\subsubsection{模型建立}

% 开篇引文
针对问题一，本节将从...角度建立...模型。具体过程如下。

% 步骤1：逐步建模，每步 textbf{步骤N：标题} + 公式推导 + 变量说明
\textbf{步骤1：...}

首先，...。设...为...，则有：
\begin{equation}
    ... = ...
\end{equation}
其中，$...$表示...。

将...代入...，可得：
\begin{equation}
    ... = ...
\end{equation}
进一步展开得：
\begin{equation}
    ... = ...
\end{equation}

\textbf{步骤2：...}

根据...的性质，可以建立以下关系：
\begin{equation}
    ... = ...
\end{equation}
其中，$...$为...，$...$为...。

由式(1)和式(2)联立可得：
\begin{equation}
    ... = ...
\end{equation}

\textbf{步骤3：...}

% ...（按需增加步骤，每步≥6个公式）...

\textbf{步骤4：...算法}

% 算法流程公式 + 参数表

\begin{longtable}{>{\centering\arraybackslash}p{0.14\textwidth}>{\raggedright\arraybackslash}p{0.30\textwidth}>{\centering\arraybackslash}p{0.12\textwidth}>{\raggedright\arraybackslash}p{0.32\textwidth}}
  \caption{算法参数设置}
  \label{tab:参数1} \\
  \toprule
  参数名 & 参数含义 & 取值 & 选取依据 \\
  \midrule
  \endfirsthead
  \caption{算法参数设置（续）} \\
  \toprule
  参数名 & 参数含义 & 取值 & 选取依据 \\
  \midrule
  \endhead
  \bottomrule
  \endfoot
  ... & ... & ... & ... \\
\end{longtable}

% 总结桥接
综上所述，本节完成了...模型的建立，得到了...（目标函数/评价体系/预测模型），接下来将基于此进行求解。
